\begin{theorembox}{Theorem (Unions and Finite Intersections of Open Sets)}
\begin{theorem}[Unions and Finite Intersections of Open Sets]
\label{thm:open-set-closure-operations}
The following statements hold in $\mathbb{R}$:
\begin{enumerate}[label=(\roman*)]
  \item The sets $\varnothing$ and $\mathbb{R}$ are open.
  \item If $\{U_\lambda\}_{\lambda\in\Lambda}$ is any family of open subsets
  of $\mathbb{R}$, then
  \[
    \bigcup_{\lambda\in\Lambda}U_\lambda
  \]
  is open.
  \item If $n\in\mathbb{N}$ and $U_1,\ldots,U_n\subseteq\mathbb{R}$ are open,
  then
  \[
    \bigcap_{k=1}^{n}U_k
  \]
  is open.
\end{enumerate}
\hyperref[prf:open-set-closure-operations]{Go to proof.}
\end{theorem}
\end{theorembox}
\begin{remark*}[Standard quantified statement]
\[
\Bigl(\forall i\in I\;U_i\text{ open}\Bigr)\Longrightarrow \Bigl(\bigcup_{i\in I} U_{i}\Bigr) \text{ open};\qquad
\Bigl(\bigwedge_{k=1}^{n} U_k\text{ open}\Bigr)\Longrightarrow \Bigl(\bigcap_{k=1}^{n} U_{k}\Bigr) \text{ open}.
\]
\end{remark*}
\begin{remark*}[Interpretation]
Openness is preserved by arbitrary unions but only finite intersections. The infinite intersection $\bigcap_n(-1/n,1/n)=\{0\}$ shows why ``finite'' cannot be dropped: a shrinking family of open intervals can close down onto a single point.
\end{remark*}
\begin{dependencies}
\begin{itemize}
  \item \hyperref[def:open-set]{Open Set in $\mathbb{R}$}
\end{itemize}
\end{dependencies}
